\subsection{\texorpdfstring{\(\bar{R}\) 中的加法与乘法}{\textbackslash bar\{R\} 中的加法与乘法}}

首先注意，函数 \(-x\) 可以连续延拓到 \(\overline{\mathbf{R}}\) ，依照公式 \(-(-\infty)=+\infty\) 与 \(-(+\infty)=-\infty\) ；如此延拓后的函数是 \(\overline{\mathbf{R}}\) 到其自身的一个同胚。

其次，考察函数 \(x + y\) 与 xy，它们定义在 \(R \times R\) 上、取值于 R 中；通过认为它们取值于拓扑空间 \(\overline{R}\) ，我们将看到它们也能在 \(\overline{R} \times \overline{R}\) 的某些点连续延拓。

就 \(x + y\) 而言，令 \(A' = ] - \infty, + \infty]\) 与 \(A'' = [-\infty, + \infty[\) 。于是有下述命题：

命题 7. 函数 \(x + y\) 可以依照公式连续延拓到集 \(\mathbf{A}' \times \mathbf{A}'\) 与 \(\mathbf{A}'' \times \mathbf{A}''\) 的每一个上，公式为

\[
\left\{ \begin{array}{l l} x + (+ \infty) = (+ \infty) + x = + \infty & (x \neq - \infty), \\ x + (- \infty) = (- \infty) + x = - \infty & (x \neq + \infty). \end{array} \right.\tag{2}
\]

例如我们证明：当 \((x,y)\) 趋于点 \((a,+\infty)\) \((a\neq-\infty)\) 而保持属于 \(R\times R\) 时，\(x+y\) 趋于 \(+\infty\) 。存在有限数 b\textless a ，而区间 {]}b, \(+\infty\) {[} 是 a 在 R 中的邻域；对任意有限 c ，关系 x\textgreater b 与 y\textgreater c-b 蕴含 \(x+y>c\) ，这表明 \(x+y\) 可以任意接近 \(+\infty\) ，只要 \((x,y)\) 充分接近 \((a,+\infty)\) 。其他情形论证类似。

相反，\(x + y\) 在点 \((-\infty, +\infty)\) 与 \((+\infty, -\infty)\) （ \(\overline{\mathbf{R}} \times \overline{\mathbf{R}}\) 的）处没有极限。因为若 \(x + y\) 有极限 \(k\) （有限或无限），当 \((x, y)\) 趋于 \((+\infty, -\infty)\) 而保持属于 \(\mathbf{R} \times \mathbf{R}\) 时，则将推出：对每个有限 \(a\) ，函数 \((x + a) - x\) 将趋于 \(k\) ，当 \(x\) 趋于 \(+\infty\) 而保持属于 \(\mathbf{R}\) 时；而这是荒谬的，因为 \((x + a) - x = a\) 且 \(a\) 是任意的。

函数 \(x + y\) 把 \(A' \times A'\) （相应地 \(A'' \times A''\) ）映入 \(A'\) （相应地 \(A''\) ）。因此它是 \(A'\) （相应地 \(A''\) ）上的一个合成法则，它延拓了 \(R\) 上的加法法则。由恒等延拓原理（第一章 § 8 第 1 目，命题 2 的推论 1），这个法则交换且结合； o 是这个法则的单位元素；而由公式 (2)，\(A'\) 的唯一非正则元素（代数，第一章 § 2 第 2 目）是 \(+\infty\) 。

若 \(x, y, z, t\) 是 \(\overline{\mathbf{R}}\) 中满足 \(x \leqslant y\) 与 \(z \leqslant t\) 的点，则只要不等式两边都有定义，就有 \(x + z \leqslant y + t\) 。

注意在 \(\overline{\mathbf{R}}\) 中，由公式 (2)，关系 \(x < y\) 蕴含 \(x + z < y + z\) 仅当 \(z\) 有限时；容易验证，只要不等式两边都有定义，关系 \(x < y\) 与 \(z < t\) 就蕴含 \(x + z < y + t\) 。

在 \(\overline{\mathbf{R}}\) 中我们令 \(x^{+} = \sup (x,0)\) ，\(x^{-} = \sup (-x,0)\) ，\(|x| = \sup (x, - x)\) ；于是 \((+\infty)^{+} = (-\infty)^{-} = +\infty\) 且 \((+\infty)^{-} = (-\infty)^{+} = 0\) ，而 \(|+\infty| = |- \infty| = +\infty\) 。和 \(x^{+} - x^{-}\) 与 \(x^{+} + x^{-}\) 对所有 \(x \in \overline{\mathbf{R}}\) 有定义，因而由恒等延拓原理分别等于 \(x\) 与 \(|x|\) 。同样，只要和 \(x + y\) 有定义，我们就有 \(|x + y| \leqslant |x| + |y|\) 。

相反要注意，§ 1 第 1 目的公式 (6) 与 (7) 对某些值 \(x\) 与 \(y\) （ \(\overline{\mathbf{R}}\) 中的）可能不再有意义；例如，若 \(x = -\infty\) 且 \(y = 0\) ，我们有 \(\sup (x, y) = 0\) ，但和 \(x + (y - x)^+\) 没有定义，因为 \((y - x)^+ = +\infty\) 。

令 \(\overline{\mathbf{R}}^{*}\) 表示 \(o\) 在 \(\overline{\mathbf{R}}\) 中的余集。于是乘法情形的命题 7 的类似命题如下：

命题 8. 函数 \(xy\) 可以依照公式连续延拓到集 \(\overline{\mathbf{R}}^* \times \overline{\mathbf{R}}^*\) 上，公式为

\[
\left\{ \begin{array}{l} x. (+ \infty) = (+ \infty). x = \left\{ \begin{array}{l l} + \infty & \text {if} \quad x > 0 \\ - \infty & \text {if} \quad x <   0 \end{array} \right. \\ x. (- \infty) = (- \infty). x = \left\{ \begin{array}{l l} - \infty & \text {if} \quad x > 0 \\ + \infty & \text {if} \quad x <   0. \end{array} \right. \end{array} \right.\tag{3}
\]

我们把证明留给读者；它与命题 7 的证明类似。

同样，我们看到 \(xy\) 在点 \((0, +\infty), (+\infty, 0), (0, -\infty), (-\infty, 0)\) （ \(\overline{\mathbf{R}} \times \overline{\mathbf{R}}\) 的）处没有极限。

函数 \(xy\) 是 \(\overline{\mathbf{R}}^*\) 上的一个合成法则，它延拓了 \(\mathbf{R}\) 上的乘法法则；这个法则结合且交换（恒等延拓原理）；它以 1 为单位元素；而 \(\overline{\mathbf{R}}^*\) 中的非正则元素是 \(+\infty\) 与 \(-\infty\) 。

若 \(x \leqslant y\) 且 \(z > 0\) ，则只要不等式两边都有定义，就有 \(xz \leqslant yz\) 。若积 \(xy\) 有定义，则 \(|x| \cdot |y|\) 也有定义，且我们有 \(|xy| = |x| \cdot |y|\) 。

最后，分配律公式

\[
x (y + z) = x y + x z\tag{4}
\]

只要出现在两边的所有运算都有定义，它由恒等延拓原理仍然有效。

\pandocbounded{\includegraphics[keepaspectratio]{images/part002_cca5eabe.jpg}}

注意可能发生 (4) 的左边有定义而右边没有定义的情形：例如考察 \(x = +\infty\) ， y = 2 与 z = -1 的情形。因此在 R 中使用分配律公式应当谨慎。

